\centerline{\bf \Huge Сложный числовой разнобой}

\begin{flushright}
        \scriptsize{--- Никогда не спрашивайте женщину о её возрасте,\\ мужчину о его зарплате и Алину зачем нужна инверсия}
\end{flushright}

\problem{1}
Пусть $a$, $b$, $c$ — произвольные натуральные числа. Докажите, что числа $a^2+b+c$, $b^2+c+a$, $c^2+a+b$ не могут одновременно быть точными квадратами.

\problem{2}
Докажите, что из любых десяти последовательных натуральных чисел можно выбрать число, взаимно простое с остальными девятью.\ Докажите, что произведение десяти последовательных натуральных чисел не может быть точным квадратом.

\problem{3}
Натуральные числа $a$, $b$ и $c$ таковы, что $ab/(a-b)=c$. Известно также, что числа $a$, $b$ и $c$ не имеют общего для них всех натурального делителя, большего 1. Докажите, что $a-b$ есть точный квадрат.

\problem{4}
Пусть $M$ — некоторое непустое (возможно, бесконечное) подмножество простых чисел. Известно, что для любого подмножества $K\subset M$ число $(\prod_{p\in K} p )- 1$ раскладывается на простые множители, принадлежащие $M$. Докажите, что тогда $M$ — это все множество простых чисел.

\problem{5}
Вася выписал в строку 100 последовательных чисел. Во второй строке под каждым числом он написал его собственный делитель. В третьей строке он под каждым числом из второй строки написал его собственный делитель и т.д., пока не получилось 2022 строк. Могут ли в каждой строке стоять последовательные числа?

\problem{6}
Пусть $p > 5$ — простое число и $S:=\{p - n^2 \colon n \in \mathbb{N}, \, n^2 < p\}$. Докажите, что во множестве $S$ найдутся два числа, одно из которых делит другое.

\problem{7}
Назовем натуральное число $N$ \textit{сбалансированным}, если или $N=1$, или $N$ раскладывается в произведение четного количества простых чисел с учетом кратности. Пусть также $P(x)=(x+a)(x+b)$ для некоторых натуральных $a$ и $b$.\ Докажите, что найдутся такие различные $a$ и $b$, что все числа $P(1)$, $P(2)$, …, $P(2022)$ сбалансированы.\ Докажите, что если $P(n)$ сбалансировано при любом натуральном $n$, то $a = b$.

\newpage

\hspace{3mm}

\problem{8}
На доске написаны числа 1, 2, 3, 4, …, 2022. Петя и Вася играют в игру. Они по очереди стирают написанные числа. Начинает Петя. Тот из мальчиков, после хода которого все оставшиеся числа (быть может, одно) будут иметь натуральный общий делитель, больший единицы, проигрывает. Если же на доске останется только единица, то игра считается завершившейся вничью. Каков будет результат при правильной игре обоих соперников?

\problem{9}
Найти все тройки натуральных чисел $(a,b,c)$, такие, что $2^a+2^b+1$ делится на $2^c+1$.

\problem{10}
Для заданного нечетного числа $n>3$ обозначим через $k$ и $\ell$ наименьшие натуральные числа такие, что числа $kn+1$ и $\ell n$ — точные квадраты. Известно, что $k$, $\ell > n/4$. Докажите, что число $n$ --- простое.